% !TeX program = pdflatex
\documentclass[runningheads]{llncs}

\usepackage[T1]{fontenc}
\usepackage{microtype}
\usepackage{booktabs}
\usepackage{tabularx}
\usepackage{longtable}
\usepackage{array}
\usepackage{enumitem}
\usepackage{xcolor}
\usepackage{hyperref}
\usepackage{xurl}
\usepackage{listings}

\hypersetup{
  colorlinks=true,
  linkcolor=blue!50!black,
  citecolor=blue!50!black,
  urlcolor=blue!60!black,
  pdftitle={A Reproducible Cross-National Dataset of Public Student Visa and Study Permit Statistics},
  pdfauthor={Vitalii Kaplan}
}

\setlist{nosep,leftmargin=*}
\setlength{\emergencystretch}{1em}
\renewcommand{\arraystretch}{1.16}
\newcolumntype{Y}{>{\raggedright\arraybackslash}X}
\newcolumntype{L}[1]{>{\raggedright\arraybackslash}p{#1}}
\newcommand{\file}[1]{\path{#1}}

\lstdefinestyle{code}{
  basicstyle=\ttfamily\small,
  backgroundcolor=\color{gray!8},
  frame=single,
  rulecolor=\color{gray!35},
  breaklines=true,
  columns=fullflexible,
  keepspaces=true,
  showstringspaces=false
}

\title{A Reproducible Cross-National Dataset of Public Student Visa and Study Permit Statistics}
\titlerunning{A Reproducible Cross-National Student-Visa Dataset}
\author{Vitalii Kaplan\orcidID{0009-0009-8181-2863}}
\authorrunning{V. Kaplan}
\institute{\email{2333275007@alu.istinye.edu.tr}}

\begin{document}

\raggedbottom
\maketitle

\begin{abstract}
Public statistics on student visas and study permits are difficult to compare because governments publish different measures, populations, time periods, origin concepts, and denominators. We present a reproducible cross-national dataset constructed from official sources for Australia, Canada, New Zealand, the United Kingdom, and the United States. The resulting SQLite artifact contains 9,510,731 aggregate measure records in a 30-field long-form schema, with overall coverage extending from the late 1990s to 2026. Applications, decisions, grants, refusals, issuances, permit holders, and reported rates remain distinct rather than being combined as nominally equivalent outcomes. Original origin labels are preserved and normalized against the United Nations M49 reference using reviewed aliases and explicit statuses for matched countries, aggregates, unknown values, and unresolved labels. The preparation workflow separates source retrieval, extraction, transformation, normalization, database loading, and validation, and can reproduce the derived artifacts from preserved source snapshots. Provenance fields connect analytical records to their packaged source files, while automated checks evaluate mappings, applicant scope, measure consistency, and database integrity. The dataset supports public-data audits, destination-specific analysis, and carefully scoped cross-national comparisons. It contains no individual application records, and differences in legal categories, reporting practices, populations, and statistical denominators limit na\"ive comparisons of approval rates.
\keywords{Educational migration \and Student visa \and Open data collection \and Historical data}
\end{abstract}

\section{Introduction}

International education research often begins with enrolment, mobility, or institution-level data \cite{abbott2016,huang2023,weber2023,bartha2018}. A separate transition occurs after admission or eligibility to study and before study begins: an applicant requests a student visa, study permit, or equivalent authorization and receives a government decision \cite{chen2023,gonnot2025}. Public statistics describing this transition can inform migration research, policy evaluation, and international-education planning. However, no assumption can be made that two official tables bearing similar labels represent the same event or denominator.

The underlying data are fragmented institutionally and technically. Immigration New Zealand publishes annual offshore or overseas decision tables on web pages \cite{inz}. The Australian Department of Home Affairs distributes large pivot-style workbooks covering lodgements, grants, and decision-based grant rates \cite{australia}. The UK Home Office publishes broad entry-clearance and sponsored-study workbooks that require category and applicant filtering \cite{uk}. Canadian records span committee-transparency pages and open-data files with different populations and origin concepts \cite{canada}. The US Department of State publishes visa issuances by nationality and visa class, but not nationality-level refusal denominators in the same table \cite{us}.

This work treats the dataset itself as a research output, following related efforts that document the construction, transformation, and limitations of reusable mobility and visa-policy data \cite{recchi2021,weber2023}. Its objective is to determine what reusable, origin-aware student-visa data can be reconstructed from these publicly available government sources while preserving their statistical meanings. This approach is consistent with the broader principle that reusable data should be findable, accessible, interoperable, and reusable, while recognizing that openness alone does not create semantic comparability \cite{fair}.

The contributions are:

\begin{itemize}
  \item a preserved inventory of official source files and their source-level schemas;
  \item reproducible extraction of HTML, CSV, ordinary XLSX, and XLSX pivot-cache records;
  \item a 30-field long schema that retains source-specific dimensions and measure definitions;
  \item destination-aware origin normalization against a preserved UN M49 reference \cite{unsd};
  \item a strict SQLite artifact that excludes non-study proxies and explicitly identifiable dependant cohorts; and
  \item executable validation and provenance audits rather than undocumented manual cleaning.
\end{itemize}

The dataset is intended for aggregate analysis. It contains no personal records, does not represent individual application histories, and must not be used as individualized visa advice or as evidence of discrimination or causality without a separately designed study.

The code and documentation are maintained in the \emph{study-visa-open-data} GitHub project \cite{repository}. Unless an absolute location is explicitly given, every relative file and directory path in this article is interpreted from the root of that project.

\section{Data Practices in Related Research}

Research on international student migration makes substantial use of public and publicly documented data. The literature links visa records with institutional information, constructs cross-national mobility and policy datasets, and supplements aggregate evidence with surveys and documentary analysis. These approaches measure different stages of mobility; applications, decisions, enrolments, student stocks, costs, and intentions must therefore remain analytically distinct.

\subsection{Administrative records and linked institutional data}

Administrative linkage studies connect visa measures to education records. Chen et al. identify international candidates in standardized-testing data, link them to enrolment and institution records, and obtain visa-refusal measures through government information requests. They evaluate coverage against independent enrolment series and acknowledge that refusal rates reflect both policy and applicant composition \cite{chen2023}. Gonnot and Lanati combine Canadian post-secondary enrolments, programme participation records, and study-permit applications and approvals. Institutional changes and missing observations are handled explicitly, while robustness tests assess the use of approvals where enrolments are unavailable \cite{gonnot2025}. Both studies demonstrate that linkage and substitution assumptions must be documented and tested.

\subsection{Constructed cross-national policy and mobility datasets}

Recchi et al. construct a global visa-cost dataset through a documented hierarchy of government and consular sources. Visa categories remain separate, currencies are normalized to a common reference point, waivers are distinguished from missing values, and difficult cases undergo repeated review. Data collected from a commercial service are retained separately after their coverage and reliability prove inconsistent with official sources \cite{recchi2021}. This work provides a strong precedent for source prioritization and for avoiding unsupported integration.

Bilateral mobility studies commonly join OECD or UNESCO student flows to World Bank indicators, CEPII geographic and historical variables, university rankings, trade measures, and national statistics. Abbott and Silles define flows by citizenship and exclude people already resident in the destination before applying gravity models \cite{abbott2016}. Huang et al. select the cross-section with the best common coverage and compare global and regional mobility structures \cite{huang2023}. Weber and Van Mol report coverage, mark imputed values, document extrapolated ranking measures, and distinguish between-country differences from within-country change \cite{weber2023}. Kler et al. combine mobility, economic, education, migration-stock, trade, and distance variables using models appropriate for zero bilateral flows \cite{kler2023}. Bartha and Gubik derive an Erasmus mobility rate from participant and student populations, while recognizing that joined public sources cover different country sets \cite{bartha2018}. Collectively, these studies emphasize population definitions, temporal alignment, coverage, derived denominators, structural zeros, and transparent treatment of imputation.

\subsection{Surveys, reviews, and descriptive policy evidence}

Surveys capture motivations that administrative counts cannot observe, but sampling determines how results may be generalized. Hercog and van de Laar combine a purposive student survey with interviews and explicitly distinguish mobility intentions from completed migration \cite{hercog2017}. Shichkin et al. combine literature searches with an expert survey, while identifying the absence of student respondents as a limitation \cite{shichkin2025}.

Documentary studies broaden interpretation but often provide less reproducible selection. Kirloskar and Inamdar synthesize journal literature, policy documents, and country reports across major destinations \cite{kirloskar2021}. Choudaha combines international mobility statistics with political and economic events to produce a qualitative periodization rather than a statistical forecast \cite{choudaha2017}. Spoonley uses aggregate migration series and policy documents to interpret changes in New Zealand's immigration system; study visas form part of that policy context, but the evidence does not represent a student-visa applicant cohort \cite{spoonley2025}.

The literature therefore confirms extensive use of public data in student-migration research. Its most transferable practices are to preserve source definitions, document transformations, distinguish missing and derived values, assess coverage, and validate results against alternative sources. These practices enable comparison without collapsing administrative events, enrolments, and mobility intentions into a single outcome.

\section{Data Sources}

\subsection{Data source references}

We used official web pages and downloadable resources published by national immigration authorities and other government organizations as the sources of data. At the time of collection, all materials were publicly available without registration, authentication, or special permission. The retrieval and extraction scripts provided in the project repository accessed and processed the sources in their published form and operated successfully at the time of data collection.

Table~\ref{tab:sources} summarizes the data sources. The complete inventory, including URLs, local filenames, formats, fields, reporting periods, and comparability notes, is machine-readable in \file{data/country-schema.json}.

\begin{table}[htbp]
\centering
\caption{Official data sources and their role in the dataset.}
\label{tab:sources}
\small
\begin{tabularx}{\textwidth}{@{}L{2.1cm}L{3.0cm}Y Y@{}}
\toprule
Destination & Publisher & Retrieved data & Principal qualification \\
\midrule
New Zealand & Immigration New Zealand \cite{inz} & Annual offshore/overseas applications, approvals, declines, and rates, 2022--2025 & The 2022 schema differs from 2023--2025; coverage is offshore/overseas only. \\
Australia & Department of Home Affairs \cite{australia} & Lodgement, grant, and grant-rate XLSX pivot workbooks, FY2005--06 through 30 April 2026 & Grant rate is grants divided by grants plus refusals, not by lodgements. Only primary applicants enter generated analytical artifacts. \\
United Kingdom & Home Office \cite{uk} & Entry-clearance applications/outcomes and sponsored-study course-level workbooks, 2005--2026 Q1 & Broad entry-clearance data require a Study-category filter. Only main applicants are retained; course-level tables cover CAS-matched main applicants. \\
Canada & Immigration, Refugees and Citizenship Canada \cite{canada} & Approval/outcome snapshots, intake and issuance tables, permit-holder files, and a French/bilingual subset, 2015--2026 with outcome subsets for 2019--April 2024 & Country of residence and citizenship are different origin concepts. Recent outcome tables have restricted populations or top-country coverage. \\
United States & Department of State \cite{us} & F1 and M1 issuances by nationality and fiscal year, FY1997--FY2024 & Issuances do not provide nationality-level application/refusal denominators; F2 and M2 dependants are excluded. \\
\bottomrule
\end{tabularx}
\end{table}

\subsection{Source acquisition and provenance}

The separate Zenodo data artifact contains the raw snapshots used for this release \cite{dataset}. Retrieval commands and official URLs are documented in \file{data/RETRIEVAL.md}, while source metadata and schemas are recorded in \file{data/country-schema.json}. The end-to-end build therefore reproduces derived files from preserved artifact inputs or from sources obtained with the documented retrieval commands. It does not assume that future requests to every remote URL will return identical bytes, because government resources and dashboard endpoints may change.

\section{Collection Methodology}

Building on this inventory, collection was organized around preserved public records and machine-readable source descriptions. Official government sources were located by publisher, statistical topic, reporting period, and availability of downloadable tables. Each retained source is identified by its official URL, publisher, access or snapshot date, local filename, format, reporting coverage, and source-specific definitions in \file{data/country-schema.json}. Retrieval commands are recorded in \file{data/RETRIEVAL.md}. When a visa-statistics source was distributed as a static CSV, XLSX, or HTML file, it was downloaded and preserved unchanged; transformations operate on repository copies rather than altering the raw input. The present release was assembled with Python~3 from source snapshots available by 19 July 2026.

\subsection{Executable build chain}

The supported build is launched from the repository root:

\begin{lstlisting}[style=code,language=bash]
python3 scripts/rebuild_database.py
\end{lstlisting}

The driver stops if any child process fails. Table~\ref{tab:pipeline} summarizes the seven documented stages. Their parameters, inputs, outputs, and observed row counts are also recorded in JSON form at \file{scripts/collect_data_script_chain.json}.

\begin{table}[htbp]
\centering
\caption{Reproducible preparation chain.}
\label{tab:pipeline}
\small
\begin{tabularx}{\textwidth}{@{}r L{4.4cm} Y@{}}
\toprule
Step & Script & Function \\
\midrule
1 & \file{extract_csv.py} & Extract New Zealand, Australian, UK, and Canadian source tables. \\
2 & \file{extract_us.py} & Extract United States F1 and M1 issuances. \\
3 & \file{combine_canada.py} & Combine Canadian outcome, intake/issuance, and permit-volume source tables without conflating their coverage. \\
4 & \file{extract_canada_recent.py} & Produce the Q-2700 French-speaking/bilingual subset for SQLite-only loading. \\
5 & \file{build_common_table.py} & Transform supported records to the common long schema and apply applicant-scope rules. \\
6 & \file{load_sqlite.py} & Apply the strict category filter, load a temporary database, create indexes, and atomically replace the final database. \\
7 & \file{validate_database.py} & Validate references, measures, mappings, provenance, reconciliation rules, and applicant scope; write the unmatched-label audit. \\
\bottomrule
\end{tabularx}
\end{table}

For a faster rebuild that reuses the source-level CSV extractions, the driver accepts the \file{--skip-extraction} option. Validation can be disabled with \file{--skip-validation}, although this is not recommended for a release build.

\subsection{Format-specific extraction}

Extraction methods follow the structure of each source format. HTML tables are parsed from preserved government pages, while ordinary CSV and XLSX tables are read according to their documented sheet and field layouts. Australian workbooks require a different method because the visible sheets are pivot reports rather than flat source tables. The extractor reads the underlying XLSX pivot-cache definitions and records directly, recovering dimensions that would be lost by copying visible cells, including lodgement channel and, for the grants workbook, the last visa held.

The extracted Australian source tables contain 2,079,645 lodgement records, 2,274,440 grant records, and 2,046,442 grant-rate/decision records before primary-applicant filtering. The UK extraction produces 249,920 broad entry-clearance application records, 402,372 broad outcome records, 18,717 sponsored-study course application records, and 18,658 course-level grant records before strict category and applicant filtering. These source-level counts should not be interpreted as numbers of individual applicants: each record is a dimensional aggregate and later may emit one or more long-form measures.

US workbook count cells are normalized to their displayed integer values. Missing or suppressed values remain missing rather than being converted to zero. Region totals and grand totals are excluded, as are dependent visa classes F2 and M2.

\section{Data Processing}

Processing converts heterogeneous extracted tables into scoped analytical records through explicit applicant and visa-category filters, country-name reconciliation, long-form transformation, and atomic loading. Incompatible measures remain separate, and unavailable values are not imputed.

\subsection{Applicant and category scope}

The strict analytical artifacts exclude explicitly identifiable dependant or mixed-applicant cohorts:

\begin{itemize}
  \item Australian source extractions preserve both values for provenance, but common outputs retain \texttt{Applicant Type = Primary};
  \item UK entry-clearance source extractions preserve \texttt{All}, \texttt{Main Applicant}, and \texttt{Dependant}, but common outputs retain \texttt{Main Applicant}; sponsored-study course tables are already main-applicant-only; and
  \item US outputs retain the primary student classes F1 and M1 and exclude F2 and M2.
\end{itemize}

Sources without a separable applicant-type field remain as published. Consequently, absence of an explicit dependant label is not proof that every aggregate from New Zealand or Canada represents primary applicants only.

The database loader additionally requires the destination-specific student or study category. It excludes non-study UK entry-clearance categories. The Canadian Q-2700 language subset is loaded only into SQLite and retains a distinct coverage label so that it is not silently merged with all-applicant Canadian outcomes.

\subsection{Source differences and structural missingness}

After applicant and category filtering, the common schema retains the union of dimensions and measures exposed by the five national statistical systems; it does not imply that every source supplies every field. An empty dimension therefore usually means that the attribute was not published or was not applicable to that source. If a measure is unavailable, the long-form transformation emits no row for it rather than creating a blank-valued row or replacing the value with zero. Derived values are produced only where the available components support the official definition.

\textbf{New Zealand.} The annual Immigration New Zealand tables describe offshore or overseas decisions by origin and calendar year, so fields for financial year, quarter, month, education sector or level, age, gender, provider state, and lodgement channel remain empty. The 2023--2025 tables report completed applications, approvals, declines, approval rate, and decline rate, but the 2022 table reports only approvals, declines, and decline rate. A 2022 completed-application total and approval rate could be calculated only if approvals and declines exhaust the source population. Because the source does not establish that withdrawals or other outcomes are absent, the present build does not make that assumption: application-volume and approval-rate rows for 2022 remain unavailable.

\textbf{Australia.} The Department of Home Affairs publishes separate monthly pivot workbooks for lodgements, grants, and decision-based grant rates. Their rich dimensions populate financial year, quarter, month, client location, lodgement channel, sector, applicant type, provider state, gender, age, and citizenship, but they do not provide a separate education-level field; the grants workbook alone exposes last visa held. Lodgements, grants, and decisions are different events and are therefore not used to fill one another's missing measures. In the decision workbook, the source provides grants, refusals, and total decisions; validation confirms that grants plus refusals equals that total, and the pipeline calculates grant rate as grants divided by grants plus refusals. The dedicated grants workbook and the grant component of the decision workbook remain distinct because their dimensional structures differ.

\textbf{United Kingdom.} UK data are quarterly and nationality-based. General entry-clearance files provide applications and decision outcomes, while sponsored-study files provide course-level applications and grants only from 2018 onward. Consequently, education level is populated for the course-level series but empty for general entry-clearance records, and fields such as applicant location, lodgement channel, age, gender, and provider state are not available. Applications and outcomes published in the same quarter do not necessarily describe the same application cohort, so the pipeline does not derive a quarterly approval rate by dividing grants by applications. Refused, withdrawn, and lapsed decisions remain separate published outcome measures.

\textbf{Canada.} Canadian coverage is assembled from several sources with different populations and origin concepts. Older outcome tables use country of residence and may report approvals, refusals, withdrawals, processed totals, and approval rates; recent snapshots may cover only leading countries, aggregates, or French-speaking/bilingual applicants. Separate Government of Canada Open Data files report permit holders or permits effective by citizenship rather than applications and decisions. These differences leave many dimensions and measures empty for particular Canadian rows and prevent one source from completing another. Source-reported approval rates are retained with their stated denominator. In the Q-2700 French-speaking/bilingual subset, processed applications can be derived as approvals plus refusals because those components define the reported decision total, but that subset remains explicitly labelled and SQLite-only rather than being treated as an all-applicant national series.

\textbf{United States.} The Department of State workbooks provide annual F1 and M1 visa issuances by nationality and fiscal year. They do not provide nationality-level applications, refusals, withdrawals, decisions, or approval rates in the same source, nor do they expose education sector, level, applicant location, age, gender, provider state, or lodgement channel. Those fields therefore remain empty or have no corresponding measure rows. An approval rate cannot be responsibly derived from issuances alone because its application or decision denominator is unavailable.

\subsection{Country-name reconciliation}

After source-specific processing, origin labels are normalized against the United Nations Statistics Division M49 methodology table. The official source was downloaded on 19 July 2026 as \file{data/UNSD — Methodology.csv} from the UNSD M49 methodology page and is used directly by the normalization scripts \cite{unsd}. The table contains 248 unique countries or areas and records the canonical name, three-digit M49 code, region, subregion, intermediate region, and ISO alpha-2 and alpha-3 codes.

Normalization first applies destination-specific or global reviewed rules from \file{data/country_aliases.csv}, then exact and normalized-form matching against the M49 canonical names. Each source label receives one of five statuses:

\begin{description}[style=nextline]
  \item[matched] an official M49 name matches directly or after orthographic normalization;
  \item[alias\_matched] a reviewed source-specific alias maps to an M49 name;
  \item[aggregate] the label denotes a total, region, or other non-country group;
  \item[unknown] the source explicitly denotes unknown or unstated origin; or
  \item[unmatched] no responsible canonical assignment is available.
\end{description}

Table~\ref{tab:normalization-status} reports the resulting distribution in the 9,510,731-row strict SQLite database. These are measure-row counts rather than counts of distinct source labels or applicants, so Australia's highly dimensional records have a large influence on the percentages.

\begin{table}[htbp]
\centering
\caption{Origin-country normalization results in the strict SQLite database.}
\label{tab:normalization-status}
\small
\begin{tabularx}{\textwidth}{@{}L{3.1cm}r r Y@{}}
\toprule
Normalization status & Rows & Share (\%) & Canonical \texttt{un\_m49\_country} \\
\midrule
\texttt{matched} & 6,806,596 & 71.568 & Assigned from the M49 reference \\
\texttt{alias\_matched} & 2,416,032 & 25.403 & Assigned through a reviewed alias \\
\texttt{aggregate} & 1,459 & 0.015 & Not assigned; label is not a country or area \\
\texttt{unknown} & 20,542 & 0.216 & Not assigned; source reports unknown origin \\
\texttt{unmatched} & 266,102 & 2.798 & Not assigned; no responsible mapping available \\
\midrule
Total & 9,510,731 & 100.000 & --- \\
\bottomrule
\end{tabularx}
\end{table}

Only \texttt{matched} and \texttt{alias\_matched} rows receive \texttt{un\_m49\_country}. Territories represented separately by UNSD are not reassigned to an administering state. The 266,102 unmatched rows represent 59 destination/source/label groups and are exported with destination, original label, source, row count, and period range to \file{log/unmatched_origin_labels.csv}.

\subsection{Long-form transformation}

Following normalization, source records are transformed into the common long form. Government files organize their statistics in different ways, and some place several measures in one record. For example, one Australian decision record can contain grants, refusals, total decisions, and grant rate. The transformation converts that source record into four rows, one for each measure. Every emitted row retains the same source dimensions, such as origin country, period, education sector, applicant location, gender, and age group. This structure allows different national measures to use one common schema without implying that similarly named measures have identical definitions.

Rows must nevertheless be aggregated within the correct source and coverage. Australian grants, for example, appear in both the dedicated grants workbook and the decision/grant-rate workbook. The two representations have different dimensional structures and are retained for different analytical purposes. Summing them together would count the same type of event twice. Queries must therefore select the intended source, measure, period, applicant scope, and relevant dimensions before calculating totals or rates.

\section{Dataset Description}

\subsection{Artifacts}

The resulting release is organized in layers (Table~\ref{tab:artifacts}), keeping raw files and extracted source tables distinguishable from harmonized records.

\begin{table}[htbp]
\centering
\caption{Principal release and build artifacts.}
\label{tab:artifacts}
\small
\begin{tabularx}{\textwidth}{@{}L{4.7cm}rY@{}}
\toprule
Artifact & Records & Scope \\
\midrule
\file{data/original/<country>/} & --- & Preserved official raw files; heterogeneous formats and schemas. \\
\file{data/csv/<country>/} & source-specific & Extracted source-level tables, including broad categories retained for provenance. \\
\file{data/student_visa_common_long.csv} & 9,707,903 & Thirty-field common long table for five destinations; retains distinguishable source categories and excludes identifiable dependant cohorts. \\
\file{data/student_visa.sqlite} & 9,510,731 & Strict student/study query database for five destinations; includes the labelled Canadian Q-2700 SQLite-only subset. \\
\file{log/unmatched_origin_labels.csv} & 59 groups & Audit of source-label and source combinations not assigned an M49 country. \\
\bottomrule
\end{tabularx}
\end{table}

\subsection{Database coverage}

Table~\ref{tab:coverage} reports database counts obtained from the rebuilt SQLite artifact. A ``record'' is a measure row, not a unique person, application, or decision. Distinct origin labels include aggregates and historical labels; distinct M49 values count normalized countries or areas only.

\begin{table}[htbp]
\centering
\caption{Coverage of the strict SQLite database.}
\label{tab:coverage}
\small
\begin{tabularx}{\textwidth}{@{}L{1.8cm}r L{1.5cm}L{2.4cm}Y@{}}
\toprule
Destination & Records & Origins (raw/M49) & Period & Included measures \\
\midrule
Australia & 9,319,202 & 236/224 & FY2005--06 to 30 Apr 2026 & Lodgements, grants, decisions, refusals, and grant rates \\
Canada & 54,729 & 283/235 & 2015--2026; outcome subsets 2019--Apr 2024 & Applications, outcomes, issuances, and permit-holder/effective counts \\
New Zealand & 2,777 & 180/178 & 2022--2025 & Applications, approvals, declines, and rates \\
United Kingdom & 122,895 & 212/203 & 2005--2026 Q1 & Study applications/outcomes and CAS-matched course grants \\
United States & 11,128 & 211/195 & FY1997--FY2024 & F1 and M1 issuances \\
\midrule
Total & 9,510,731 & --- & --- & --- \\
\bottomrule
\end{tabularx}
\end{table}

Australia accounts for 98.0\% of database records. This reflects the dimensional granularity of its pivot caches and the emission of multiple measures from decision rows, not a statement that 98.0\% of international student visa applicants apply to Australia. Analyses should group or sample at an explicitly defined statistical level rather than treating every long-form row as exchangeable.

\subsection{Common schema}

All SQLite columns are stored as text to preserve source formatting and allow uniform bulk loading. Consumers should cast \texttt{measure\_value} according to \texttt{measure\_unit}; validation ensures counts are non-negative integer strings and percentages lie between 0 and 100. The 30 fields are grouped in Table~\ref{tab:schema}.

\begin{longtable}{@{}L{2.4cm}L{4.1cm}L{5.3cm}@{}}
\caption{Common long-table data dictionary.}\label{tab:schema}\\
\toprule
Group & Fields & Interpretation \\
\midrule
\endfirsthead
\toprule
Group & Fields & Interpretation \\
\midrule
\endhead
Identity and origin & \texttt{destination\_country}, \texttt{origin\_country}, \texttt{un\_m49\_country}, \texttt{normalization\_status}, \texttt{origin\_country\_type} & Destination; source label; normalized M49 name; mapping status; and whether origin means citizenship, nationality, residence, or another published concept. \\
Time & \texttt{period\_label}, \texttt{period\_type}, \texttt{calendar\_year}, \texttt{financial\_year}, \texttt{quarter}, \texttt{month} & Original period representation plus parsed components where meaningful. Blank fields are explicit non-applicability or unavailable detail. \\
Visa and applicant & \texttt{visa\_category}, \texttt{visa\_category\_source}, \texttt{applicant\_location}, \texttt{lodgement\_channel}, \texttt{applicant\_type} & Harmonized and original legal/category labels, location, application channel, and applicant population. \\
Education and demographics & \texttt{education\_sector}, \texttt{education\_level}, \texttt{gender}, \texttt{age\_group}, \texttt{provider\_state}, \texttt{last\_visa\_held} & Optional dimensions retained when a source exposes them; not uniformly available across destinations. \\
Measure & \texttt{measure\_type}, \texttt{measure\_value}, \texttt{measure\_unit}, \texttt{rate\_denominator\_type} & Published or computed statistic, value, count/percent unit, and explicit rate denominator when known. \\
Coverage and provenance & \texttt{coverage\_type}, \texttt{source\_snapshot}, \texttt{source\_file}, \texttt{source\_notes} & Population/subset label, source release date, repository-relative lineage, and qualifications. \\
\bottomrule
\end{longtable}

\subsection{Indexes and query access}

SQLite indexes support destination, raw origin, origin type, period, measure, visa category, normalized M49 country, normalization status, and the composite destination--origin--measure path. A typical scoped query is:

\begin{lstlisting}[style=code,language=SQL]
SELECT financial_year, origin_country, measure_type, measure_value
FROM student_visa_common_long
WHERE destination_country = 'Australia'
  AND applicant_type = 'Primary'
  AND education_sector = 'Higher Education Sector'
  AND measure_type IN ('decisions', 'grants', 'refusals');
\end{lstlisting}

The query intentionally selects the dimensional and measure fields required for interpretation. Summing \texttt{measure\_value} without constraining source, measure, period, and relevant dimensions may double count published aggregates.

\subsection{Coverage and update frequency}

Coverage is historical but source-dependent: the earliest included US fiscal year is 1997, while other destination series begin later and may end in a quarter, month, partial year, or source-specific snapshot. There is no single update frequency for the unified dataset. Australian program workbooks are released periodically with monthly cut-offs; UK and New Zealand tables follow their official statistical releases; and several Canadian components are annual, periodic, or one-off transparency extracts. The database is therefore versioned as a snapshot and updated by rerunning the complete build against newly preserved source releases, not by assuming a universal calendar schedule.

\section{Technical Validation}

Validation is part of the default build rather than a post hoc manual review. Verification covers reference-schema integrity, dimensional-key uniqueness, decision reconciliation, numeric domains, provenance completeness, normalization consistency across artifacts, and exclusion of explicitly identifiable dependant cohorts. The release passed the checks detailed below.

\subsection{Reference and normalization checks}

The retained M49 table must contain the required schema, 248 unique non-empty country or area names, and valid three-digit M49 codes. Alias tests verify reviewed Australian label mappings, distinguish non-country labels, and ensure separately listed territories remain distinct. Normalization results must agree across the common CSV, the Canadian SQLite-only subset, and SQLite. Exact official M49 names may not remain unmatched, and aggregate, unknown, and unmatched labels may not receive a canonical country.

The normalization distribution is reported in Table~\ref{tab:normalization-status}. The 59 unmatched destination/source/label groups chiefly involve labels absent from the supplied current M49 list, including historical countries or areas, subnational areas, Taiwan, and Kosovo; the pipeline leaves them explicit rather than silently imposing a contested or historically incorrect mapping.

\subsection{Source-structure and arithmetic checks}

For the audited Australian Higher Education/Primary slice, records must be unique on all pivot-cache dimensions. This check guards against the earlier omission of lodgement channel and last-visa-held fields. Every Australian grant-rate record must satisfy

\begin{equation}
  \mathrm{grants} + \mathrm{refusals} = \mathrm{decisions}.
\end{equation}

Count values in the Australian source tables and the final database must be non-negative integers. Percentage values must lie in the closed interval $[0,100]$. Every SQLite row must contain a repository-relative source-file value.

\subsection{Scope checks}

The common CSV and SQLite are scanned for Australian \texttt{Secondary}, UK \texttt{Dependant} and \texttt{All}, generic \texttt{dependent}, and US F2/M2 records. The validated count of such explicitly excluded records in SQLite is zero. The rebuilt SQLite table contains 9,510,731 rows after strict filtering and insertion of the explicitly labelled Canadian SQLite-only subset.

Four unit tests cover reviewed aliases, non-country handling, territory preservation, and equality between generated and source M49 names and counts. The full build and all validation checks completed successfully for the 19 July 2026 snapshot.

\section{Usage Notes}

The heterogeneity preserved by the preceding workflow imposes several constraints on responsible analysis.

\subsection{No universal approval-rate denominator}

The dataset deliberately avoids constructing a universal \texttt{approval\_rate}. New Zealand's recent tables relate completed applications to approvals and declines. Australia's grant-rate workbook defines the denominator as grants plus refusals. UK applications and decisions can occur in different quarters. Canadian tables include processed outcomes, permit issuances, permit holders, and subsets with different coverage. The US table contains issuances without nationality-level application or refusal totals. These measures must not be substituted for one another.

\subsection{Origins are not semantically identical}

Depending on the source, origin is citizenship, nationality, country of residence, or an aggregate region. For Canada in particular, approval/outcome tables use country of residence while permit-holder files use citizenship. The \texttt{origin\_country\_type} field must be included in joins and comparisons. A common M49 label reconciles spelling; it does not make residence equivalent to citizenship.

\subsection{Populations and coverage vary}

New Zealand data cover offshore or overseas applications. Australia is limited in generated artifacts to primary applicants, but its three workbook families have different event definitions and dimensions. UK course-level records include main applicants matched to Confirmation of Acceptance for Studies data and should not be treated as a complete copy of entry-clearance totals. Recent Canadian tables may cover top countries, aggregates, or French-speaking/bilingual applicants rather than all applicants. US records cover F1/M1 visas issued, not all study authorizations or decisions. Applicant-type fields are absent from some sources, so primary-only coverage cannot always be verified.

\subsection{Time bases and policy changes}

The collection mixes calendar years, calendar quarters, fiscal years, financial-year months, partial years, and source-defined reporting periods. Visa categories and policy rules may change within the historical coverage. Researchers should use the parsed time fields together with \texttt{period\_label}, \texttt{source\_snapshot}, and source notes, and should not assume that a category name is legally stable over two decades.

\subsection{Long-form duplication and weighting}

Rows are published measures at dimensional keys, not statistically independent observations. One source record can emit several measure rows, and similar concepts can appear in separate official workbooks. Australia dominates the row count because it exposes more dimensions. Model training, descriptive summaries, and visualizations therefore require a defined unit of analysis and explicit aggregation rules; raw row weighting would mostly measure publication structure.

\subsection{Missingness, suppression, and normalization}

Blank fields are not imputed. A blank can mean unavailable, suppressed, not applicable, or absent from a source schema; the source documentation should be consulted before assigning a missing-data mechanism. Aggregate, unknown, and unmatched origin statuses remain distinct. Users may propose additional aliases, but changes should be reviewed and recorded rather than implemented through ad hoc string replacement.

\subsection{Appropriate and inappropriate uses}

Appropriate uses include audits of public-data availability, destination-specific descriptive statistics, reproducibility studies, and comparisons restricted to demonstrably compatible measures and populations. The resource is not suitable by itself for individualized probability estimates, causal attribution, or claims about discriminatory treatment. Such research would require careful population definitions, exposure denominators, policy controls, and an ethical and legal design beyond this data descriptor.

\section{Data and code availability}

The primary query-ready release file is \file{data/student_visa.sqlite}. The database, common CSV, raw sources, source-level CSVs, and validation audit are distributed in the Zenodo data artifact \cite{dataset}, rather than committed to the GitHub code repository. The code repository contains the scripts, machine-readable source inventory at \file{data/country-schema.json}, retrieval documentation at \file{data/RETRIEVAL.md}, and executable workflow description at \file{scripts/collect_data_script_chain.json} \cite{repository}. The data files can be obtained from the artifact or recreated with the documented retrieval and build scripts.

The cited Zenodo record identifies this release as version 1.0 under DOI \url{https://doi.org/10.5281/zenodo.21446641}. The code repository is distributed under the MIT License, while the data artifact is licensed under Creative Commons Attribution 4.0. Licences and attribution requirements of the underlying government sources remain independently applicable.

\section{Reproducibility checklist}

A release can be independently checked from the repository root as follows:

\begin{lstlisting}[style=code,language=bash]
python3 -m json.tool data/country-schema.json
python3 -m json.tool scripts/collect_data_script_chain.json
python3 scripts/rebuild_database.py
python3 -m unittest discover -s tests -v
sqlite3 -readonly data/student_visa.sqlite \
  "SELECT COUNT(*) FROM student_visa_common_long;"
\end{lstlisting}

The expected final query result for this snapshot is \texttt{9510731}. Reproduction requires Python 3, the preserved files under \file{data/original/} obtained from the separate artifact or retrieved using \file{data/RETRIEVAL.md}, and the official reference at \file{data/UNSD — Methodology.csv}. The documented build scripts use the Python standard library.

\section{Conclusion}

This dataset demonstrates that a substantial origin-aware resource can be reconstructed from public student-visa and study-permit statistics, but also that nominally similar government measures are not automatically comparable. The central methodological result is therefore both an artifact and a boundary: 9.5 million strictly scoped measure records can be queried through a common schema, while source-specific legal categories, populations, origin concepts, time bases, and denominators remain visible. Preserving those distinctions makes the resource more useful for responsible comparative research than a smaller but falsely uniform approval-rate table.

\begin{credits}
\subsubsection{\ackname}
The dataset depends on public statistical releases produced by Immigration New Zealand; the Australian Department of Home Affairs; the UK Home Office; Immigration, Refugees and Citizenship Canada; the US Department of State; and the United Nations Statistics Division.

\subsubsection{\discintname}
The author declares no competing interests.
\end{credits}

\input{references.bbl}

\end{document}
